% GENERATED FILE. Edit canonical lessons, then run npm run generate:books.
% canonical-source-hash: fnv1a64:0673012f96e8e85b
% canonical-lessons: GE-C377-hochschule, GE-C377-vorlesung, GE-C377-dozent, GE-C377-studiengang, GE-C377-praktikum, GE-R377-review-370-377

\chapter{An der Hochschule}
\label{ch:ge-an-der-hochschule}
\hlchaptermodality{377}

\begin{chapteropening}
\textbf{By the end of this chapter:} \emph{I can talk about higher education: lectures, lecturers, a degree programme and an internship.}

Complete the last lesson of chapter 377: I can talk about higher education: lectures, lecturers, a degree programme and an internship.
\end{chapteropening}

\section[die Hochschule]{die Hochschule — the university, the college}

\label{lesson:GE-C377-hochschule}



\begin{quote}
Before the new one: say the German for concentration, then the German for the revision.
\end{quote}

\subsection*{You'll want to know: die Hochschule}
\textbf{die Hochschule} — \textquotedblleft{}the university, the college\textquotedblright{}.

\begin{quote}
Meine Schwester studiert an einer Hochschule für Musik.
\end{quote}

\begin{quote}
\emph{My sister studies at a college of music.}
\end{quote}

Plural: \emph{die Hochschulen}. A general word for universities and colleges; \emph{die Fachhochschule} is a university of applied sciences.

\begin{cousinweb}
\emph{Hochschule} joins \emph{hoch} \textquotedblleft{}high\textquotedblright{}, Old High German \emph{hōh} (EtymDB), to \emph{Schule} \textquotedblleft{}school\textquotedblright{}, from Old High German \emph{scuola}, Latin \emph{schola} and Greek \emph{schol\'{\={e}}} \textquotedblleft{}leisure, learned talk\textquotedblright{} (EtymDB; Lewis \& Short).
\end{cousinweb}

\begin{grammarlens}[title={What you've built}]
One more word for telling, explaining and describing at B1.
\end{grammarlens}

\subsection*{Your turn}
\begin{itemize}
\raggedright
  \item \emph{Say it:} \emph{die Hochschule}
  \item \emph{Say it:} \emph{die Hochschule}, once more
  \item \emph{Say it:} a sentence of your own with \emph{die Hochschule}
  \item \emph{Recall:} say the German for concentration, then the German for the revision, then say \emph{die Hochschule} again
\end{itemize}

\begin{tcolorbox}[breakable,colback=teal!4,colframe=teal!35!black,title={Before you move on}]
What does die Hochschule mean? (\textquotedblleft{}The university, the college\textquotedblright{}.) Say it once more.
\end{tcolorbox}
\section[die Vorlesung]{die Vorlesung — the lecture}

\label{lesson:GE-C377-vorlesung}



\begin{quote}
Before the new one: say the German for the revision, then the German for the university.
\end{quote}

\subsection*{You'll want to know: die Vorlesung}
\textbf{die Vorlesung} — \textquotedblleft{}the lecture\textquotedblright{}.

\begin{quote}
Die Vorlesung beginnt um neun Uhr im großen Saal.
\end{quote}

\begin{quote}
\emph{The lecture starts at nine in the big hall.}
\end{quote}

Plural: \emph{die Vorlesungen}. From \emph{vorlesen} — to read aloud: lecturers once read from a book.

\begin{cousinweb}
\emph{Vorlesung} is the noun of \emph{vorlesen} \textquotedblleft{}to read aloud\textquotedblright{}, built on \emph{lesen} \textquotedblleft{}to read\textquotedblright{}, inherited from Old High German \emph{lesan} and Proto-Germanic \emph{lesaną} \textquotedblleft{}to gather\textquotedblright{} (EtymDB).
\end{cousinweb}

\begin{grammarlens}[title={What you've built}]
One more word for telling, explaining and describing at B1.
\end{grammarlens}

\subsection*{Your turn}
\begin{itemize}
\raggedright
  \item \emph{Say it:} \emph{die Vorlesung}
  \item \emph{Say it:} \emph{die Vorlesung}, once more
  \item \emph{Say it:} a sentence of your own with \emph{die Vorlesung}
  \item \emph{Recall:} say the German for the revision, then the German for the university, then say \emph{die Vorlesung} again
\end{itemize}

\begin{tcolorbox}[breakable,colback=teal!4,colframe=teal!35!black,title={Before you move on}]
What does die Vorlesung mean? (\textquotedblleft{}The lecture\textquotedblright{}.) Say it once more.
\end{tcolorbox}
\section[der Dozent]{der Dozent — the lecturer}

\label{lesson:GE-C377-dozent}



\begin{quote}
Before the new one: say the German for the university, then the German for the lecture.
\end{quote}

\subsection*{You'll want to know: der Dozent}
\textbf{der Dozent} — \textquotedblleft{}the lecturer\textquotedblright{}.

\begin{quote}
Der Dozent erklärt das Thema sehr gut.
\end{quote}

\begin{quote}
\emph{The lecturer explains the topic very well.}
\end{quote}

A weak noun: \emph{den Dozenten}, plural \emph{die Dozenten}. A woman is \emph{die Dozentin}.

\begin{cousinweb}
\emph{Dozent} is a learned borrowing from Latin \emph{docēns} \textquotedblleft{}teaching\textquotedblright{} (EtymDB), the present participle of \emph{docēre} \textquotedblleft{}to teach, to instruct\textquotedblright{} (Lewis \& Short).
\end{cousinweb}

\begin{grammarlens}[title={What you've built}]
One more word for telling, explaining and describing at B1.
\end{grammarlens}

\subsection*{Your turn}
\begin{itemize}
\raggedright
  \item \emph{Say it:} \emph{der Dozent}
  \item \emph{Say it:} \emph{der Dozent}, once more
  \item \emph{Say it:} a sentence of your own with \emph{der Dozent}
  \item \emph{Recall:} say the German for the university, then the German for the lecture, then say \emph{der Dozent} again
\end{itemize}

\begin{tcolorbox}[breakable,colback=teal!4,colframe=teal!35!black,title={Before you move on}]
What does der Dozent mean? (\textquotedblleft{}The lecturer\textquotedblright{}.) Say it once more.
\end{tcolorbox}
\section[der Studiengang]{der Studiengang — the degree programme, the programme of study}

\label{lesson:GE-C377-studiengang}



\begin{quote}
Before the new one: say the German for the lecture, then the German for the lecturer.
\end{quote}

\subsection*{You'll want to know: der Studiengang}
\textbf{der Studiengang} — \textquotedblleft{}the degree programme, the programme of study\textquotedblright{}.

\begin{quote}
Mein Studiengang dauert sechs Semester.
\end{quote}

\begin{quote}
\emph{My degree programme lasts six semesters.}
\end{quote}

Plural: \emph{die Studiengänge}. \emph{Einen Studiengang wählen} — to choose a course.

\begin{cousinweb}
\emph{Studiengang} joins \emph{Studium} \textquotedblleft{}study\textquotedblright{}, from Latin \emph{studium} \textquotedblleft{}zeal, study\textquotedblright{} (Lewis \& Short), to \emph{Gang} \textquotedblleft{}course, walk\textquotedblright{}, Old High German \emph{gang} (EtymDB): the path your studies follow.
\end{cousinweb}

\begin{grammarlens}[title={What you've built}]
One more word for telling, explaining and describing at B1.
\end{grammarlens}

\subsection*{Your turn}
\begin{itemize}
\raggedright
  \item \emph{Say it:} \emph{der Studiengang}
  \item \emph{Say it:} \emph{der Studiengang}, once more
  \item \emph{Say it:} a sentence of your own with \emph{der Studiengang}
  \item \emph{Recall:} say the German for the lecture, then the German for the lecturer, then say \emph{der Studiengang} again
\end{itemize}

\begin{tcolorbox}[breakable,colback=teal!4,colframe=teal!35!black,title={Before you move on}]
What does der Studiengang mean? (\textquotedblleft{}The degree programme, the programme of study\textquotedblright{}.) Say it once more.
\end{tcolorbox}
\section[das Praktikum]{das Praktikum — the internship, the work placement}

\label{lesson:GE-C377-praktikum}



\begin{quote}
Before the new one: say the German for the lecturer, then the German for the degree programme.
\end{quote}

\subsection*{You'll want to know: das Praktikum}
\textbf{das Praktikum} — \textquotedblleft{}the internship, the work placement\textquotedblright{}.

\begin{quote}
Im Sommer mache ich ein Praktikum bei einer Zeitung.
\end{quote}

\begin{quote}
\emph{In the summer I am doing an internship at a newspaper.}
\end{quote}

Plural: \emph{die Praktika}. \emph{Der Praktikant, die Praktikantin} — the intern.

\begin{cousinweb}
\emph{Praktikum} is a learned formation on Medieval Latin \emph{practicus} \textquotedblleft{}active, practical\textquotedblright{}, from Greek \emph{praktikós} \textquotedblleft{}fit for action\textquotedblright{}, from \emph{prássein} \textquotedblleft{}to do\textquotedblright{}. EtymDB has no entry for the German word.
\end{cousinweb}

\begin{grammarlens}[title={What you've built}]
One more word for telling, explaining and describing at B1.
\end{grammarlens}

\subsection*{Your turn}
\begin{itemize}
\raggedright
  \item \emph{Say it:} \emph{das Praktikum}
  \item \emph{Say it:} \emph{das Praktikum}, once more
  \item \emph{Say it:} a sentence of your own with \emph{das Praktikum}
  \item \emph{Recall:} say the German for the lecturer, then the German for the degree programme, then say \emph{das Praktikum} again
\end{itemize}

\begin{tcolorbox}[breakable,colback=teal!4,colframe=teal!35!black,title={Before you move on}]
What does das Praktikum mean? (\textquotedblleft{}The internship, the work placement\textquotedblright{}.) Say it once more.
\end{tcolorbox}
\section[(dialogue)]{Words from chapters 370–377 — a review}

\label{lesson:GE-R377-review-370-377}



\begin{quote}
Say the German for to work out and to go wrong.
\end{quote}

\subsection*{Your turn}
Say these aloud:

\begin{itemize}
\raggedright
  \item to work out, to go wrong, to happen, to be true, to turn out
  \item to accept, the acceptance, to call off, the rejection, the confirmation
  \item to agree on, the arrangement, to arrange, the schedule, to change your plans
  \item unforeseen, spontaneous, to take into account, the forecast, to allow for
  \item to make up your mind, the resolution, the way out, flexible, the planning
  \item the written exam, to fail an exam, oral, written, the school report
  \item to cram, by heart, to concentrate, concentration, the revision
  \item the university, the lecture, the lecturer, the degree programme, the internship
\end{itemize}

\begin{tcolorbox}[breakable,colback=teal!4,colframe=teal!35!black,title={Before you move on}]
To work out? (\textbf{klappen}.) To go wrong? (\textbf{misslingen}.) To happen? (\textbf{eintreten}.) To be true? (\textbf{zutreffen}.) To turn out? (\textbf{herausstellen}.) To accept? (\textbf{zusagen}.) The acceptance? (\textbf{die Zusage}.) To call off? (\textbf{absagen}.) The rejection? (\textbf{die Absage}.) The confirmation? (\textbf{die Bestätigung}.) To agree on? (\textbf{vereinbaren}.) The arrangement? (\textbf{die Vereinbarung}.) To arrange? (\textbf{ausmachen}.) The schedule? (\textbf{der Zeitplan}.) To change your plans? (\textbf{umplanen}.) Unforeseen? (\textbf{unvorhergesehen}.) Spontaneous? (\textbf{spontan}.) To take into account? (\textbf{berücksichtigen}.) The forecast? (\textbf{die Prognose}.) To allow for? (\textbf{einplanen}.) To make up your mind? (\textbf{entschließen}.) The resolution? (\textbf{der Entschluss}.) The way out? (\textbf{der Ausweg}.) Flexible? (\textbf{flexibel}.) The planning? (\textbf{die Planung}.) The written exam? (\textbf{die Klausur}.) To fail an exam? (\textbf{durchfallen}.) Oral? (\textbf{mündlich}.) Written? (\textbf{schriftlich}.) The school report? (\textbf{das Zeugnis}.) To cram? (\textbf{pauken}.) By heart? (\textbf{auswendig}.) To concentrate? (\textbf{konzentrieren}.) Concentration? (\textbf{die Konzentration}.) The revision? (\textbf{die Wiederholung}.) The university? (\textbf{die Hochschule}.) The lecture? (\textbf{die Vorlesung}.) The lecturer? (\textbf{der Dozent}.) The degree programme? (\textbf{der Studiengang}.) The internship? (\textbf{das Praktikum}.)
\end{tcolorbox}
